% Generado desde AnteProyecto-Final.pdf — párrafos completos

\capitulo{Resumen}

El presente documento describe el diseño y la validación empirica de un protocolo distribuido de control topológico para redes multisalto, enfocado en dispositivos móviles Android que utilizan interfaces Wi-Fi convencionales. El protocolo permite a los terminales descubrirse mutuamente, formar enlaces y organizar de forma autónoma una estructura de enrutamiento en entornos carentes de infraestructura centralizada (tales como routers o puntos de acceso fisicos). Para asegurar la continuidad lógica de la comunicación, se definen mecanismos de reconfiguración que detectan la desconexión abrupta de nodos intermedios ---un evento caracteristico de escenarios post-desastre--- y enrutan nuevamente el flujo de información, limitando los ciclos topológicos a través de tablas de estado locales guiadas por métricas de salto. La validación se llevará a cabo a nivel de capa de aplicación, probando su eficacia de resiliencia inter-nodo sin requerir modificaciones profundas a nivel de sistema operativo (root). Palabras clave: Protocolo distribuido, Wi-Fi convencional, Redes multisalto, Topologia dinámica, Android, Comunicación post-desastre.

\newpage
\capitulo{Abstract}

This document details the design and empirical validation of a distributed topological control protocol for multi-hop networks, aimed at Android mobile devices using conventional Wi-Fi interfaces. The protocol enables terminals to mutually discover each other, establish links, and autonomously organize a routing structure in environments completely lacking centralized infrastructure (such as standard physical routers or access points). To ensure the logical continuity of communication, reconfiguration mechanisms are defined to detect the abrupt disconnection of intermediate nodes---an event characteristic of post-disaster scenarios---and dynamically reroute the information flow, preventing topological cycles through local state tables guided by hop metrics. The validation will be conducted at the application layer, testing its inter-node resilience without requiring extensive operating system modifications (root access). Keywords: Distributed protocol, Conventional Wi-Fi, Multi-hop networks, Dynamic topology, Android, Post-disaster communication.

\newpage
\capitulo{Introducción}

Las redes de telecomunicaciones modernas descansan sobre arquitecturas estrictamente centralizadas, una dependencia estructural que garantiza altas velocidades y estabilidad cuando las operadoras funcionan en condiciones nominales, pero que se convierte en una falla sistémica ante contextos de disrupción, tales como desastres naturales o apagones sostenidos. En estos escenarios, el colapso de las antenas, los puntos de acceso o el tendido de fibra óptica deja a múltiples comunidades aisladas, paradójicamente, a pesar de que los dispositivos inteligentes en sus manos poseen potentes interfaces radiales capaces de emitir y recibir información a nivel local. Este proyecto aborda la formulación algoritmica y estructural de una solución orientada a la creación de redes de emergencia (redes ad hoc multisalto), eliminando por definición la necesidad de coordinadores centrales fijos. El enfoque tecnológico aprovecha el espectro Wi-Fi nativo e inalterado en dispositivos con el sistema operativo Android, sorteando limitantes de acceso a las capas MAC que históricamente han marginado a las implementaciones de redes en malla (Mesh) hacia infraestructuras costosas, hardwares estandarizados por consorcios o móviles modificados en su núcleo (rooteados). El desarrollo se fundamenta en un enfoque metodológico iterativo que permite la validación progresiva de los algoritmos de enrutamiento y reconfiguración, buscando consolidar un protocolo que garantice la conectividad lógica en redes móviles descentralizadas. El presente documento se organiza en cinco capitulos: el Capitulo 1 define el planteamiento del problema, los objetivos y el alcance del proyecto; el Capitulo 2 establece el marco de referencia teórico, legal y el estado del arte; el Capitulo 3 detalla el componente metodológico basado en incrementos funcionales; finalmente, los Capitulos 4 y 5 presentan el cronograma de ejecución y el presupuesto requerido para la materialización de la investigación.

\newpage
\capitulo{Antecedentes}

\seccion{Descripción del problema}

En la actualidad, las redes de telecomunicaciones dependen principalmente de infraestructura centralizada, como torres celulares, routers o puntos de acceso que coordinan la comunicación entre los dispositivos. Este modelo permite gestionar la conectividad de forma eficiente en condiciones normales, sin embargo, produce una dependencia directa: cuando dicha infraestructura falla, se congestiona o queda fuera de servicio, la comunicación entre los usuarios puede verse interrumpida de manera parcial o total (Tanenbaum \& Wetherall, 2011). Esta situación se vuelve especialmente relevante en escenarios como desastres naturales, fallas generalizadas de red o contextos donde la infraestructura de comunicaciones no está disponible. En estos casos, aunque los dispositivos móviles continúan funcionando y mantienen capacidades de conectividad inalámbrica, los usuarios pueden quedar aislados debido a la ausencia de sistemas que permitan coordinar la comunicación de forma directa entre los dispositivos cercanos (Akyildiz et al., 2005). A pesar de que tecnologias como Wi-Fi están ampliamente disponibles en teléfonos móviles y otros dispositivos electrónicos de uso cotidiano, su funcionamiento convencional está orientado principalmente a un modelo basado en infraestructura, donde un punto de acceso central gestiona el tráfico y la organización de la red (IEEE, 2020). Como resultado, el alcance efectivo de comunicación queda limitado al rango individual del punto central del Wi-Fi, que normalmente se encuentra aproximadamente entre 10 y 100 metros, dependiendo de las condiciones del entorno, la potencia de transmisión y las caracteristicas del dispositivo (Gupta \& Kumar, 2000). Esto conlleva a que dispositivos fisicamente cercanos pierdan su capacidad de interconexión ante cualquier fallo, saturación o limitación en la cobertura del nodo central, quedando aislados a pesar de encontrarse dentro del rango radial de otros dispositivos. Las redes multisalto presentan una alternativa que permite que la comunicación se mantenga activa incluso si uno o varios nodos fallan, aumentando la resiliencia del sistema. Además, este enfoque resulta especialmente relevante en escenarios criticos, como desastres naturales, eventos masivos o zonas rurales sin infraestructura estable, donde la comunicación inmediata y confiable es esencial (Akyildiz et al., 2005). Sin embargo, aunque existen protocolos multisalto, la mayoria requiere hardware especializado y configuraciones manuales de caracter complejo, especialmente a la hora de añadir nuevos nodos a la red, lo que limita su aplicabilidad en dispositivos de uso cotidiano (IEEE, 2020). En consecuencia, surge la necesidad de desarrollar un protocolo distribuido capaz de organizar dispositivos Wi-Fi convencionales en redes multisalto adaptativas, asegurando que la comunicación se mantenga aún en entornos dinamicos, extendiendo el alcance de comunicación y eliminando la dependencia de infraestructura centralizada.

\seccion{Formulación del problema}

Las redes inalámbricas convencionales dependen de infraestructura centralizada que, cuando falla o no está disponible, impide la comunicación entre dispositivos cercanos. Aunque los estándares IEEE 802.11 contemplan modos ad hoc y mesh (IEEE, 2020), las soluciones existentes requieren hardware especializado o modificaciones en las capas inferiores de la pila de protocolos, lo que las hace inviables en dispositivos móviles de uso cotidiano. Esto evidencia la necesidad de mecanismos de control distribuido que permitan a los dispositivos organizarse autónomamente, establecer relaciones topológicas coherentes y adaptarse a cambios en la conectividad operando exclusivamente sobre las interfaces estándar disponibles. A partir de esta necesidad se plantea la siguiente pregunta de investigación: ¿Cómo diseñar un protocolo distribuido que permita formar y mantener una topologia multisalto mediante interfaces Wi-Fi estándar, asegurando estabilidad, prevención de ciclos y reconfiguración ante cambios dinámicos?

\newpage
\capitulo{Justificación}

El estudio de mecanismos alternativos de comunicación resulta relevante en contextos donde la infraestructura tradicional de telecomunicaciones no se encuentra disponible o presenta fallas.

En situaciones como desastres naturales, interrupciones del servicio o saturación de las redes, los sistemas centralizados pueden dejar de operar temporalmente, lo que provoca que los usuarios pierdan la posibilidad de comunicarse entre si (Akyildiz et al., 2005). Esta situación se intensifica en los escenarios de desastres naturales, donde la comunicación pasa a ser un elemento vital en la busqueda de sobrevivientes y en la coordinación de las labores de rescate. Ante este tipo de escenarios, investigar formas de comunicación que no dependan de infraestructura fija puede aportar alternativas para mantener el intercambio de información entre dispositivos cercanos. Desde una perspectiva social, la posibilidad de que los dispositivos móviles puedan conectarse directamente entre si y extender la comunicación mediante saltos entre nodos puede resultar útil en entornos donde la conectividad es limitada o inexistente, como en zonas rurales, donde la infraestructura de telecomunicaciones puede ser escasa o costosa, de igual manera que en eventos masivos, donde un solo nodo central puede verse saturado por la cantidad de usuarios simultáneos (Appiahene et al., 2019). El uso de tecnologias que ya están integradas en la mayoria de los dispositivos, como el Wi-Fi, permite explorar soluciones que no requieren hardware adicional ni modificaciones complejas en la infraestructura existente, ni desde el lado de los dispositivos de los usuarios finales, como también de la infraestructura de telecomunicaciones móvil o fija. En el ámbito de la ingenieria de sistemas, el análisis de protocolos distribuidos aplicados a redes inalámbricas representa un campo de investigación importante, particularmente en temas relacionados con la autoorganización de nodos, el control de topologias dinámicas y la continuidad de la comunicación en redes con conectividad variable (Gupta \& Kumar, 2000; Karl \& Willig, 2005). Estudiar cómo estos mecanismos pueden implementarse sobre interfaces Wi-Fi estándar permite comprender mejor las posibilidades y limitaciones de este tipo de redes, especialmente ante el uso masivo moderno de la conectividad inalámbrica, misma que induce entropia en la topologia de la red. En conjunto, el desarrollo de este trabajo se enmarca dentro del proyecto de investigación denominado (Una plataforma escalable de comunicación para escenarios post-desastre utilizando tecnologias mesh y microcomputadores para redes ad hoc con cobertura ampliada). La investigación aporta al análisis de mecanismos de comunicación descentralizados que puedan operar utilizando tecnologias ampliamente disponibles, como el Wi-Fi en dispositivos móviles. Asimismo, permite generar una base conceptual y experimental que pueda servir como referencia para futuros desarrollos orientados a mejorar la resiliencia de las comunicaciones en escenarios de desastres naturales, donde la infraestructura de telecomunicaciones convencional no se encuentra disponible.

\newpage
\capitulo{Objetivos}

\seccion{Objetivo general}

Desarrollar un protocolo distribuido de control y reconfiguración de redes multisalto mediante interfaces Wi-Fi estándar, que permita a los nodos mantener comunicación estable, prevenir ciclos y adaptarse dinámicamente a cambios en la topologia, mediante un diseño experimental en dispositivos Android.

\seccion{Objetivos específicos}

Diseñar el protocolo distribuido de control y reconfiguración topológica, definiendo mecanismos de descubrimiento de nodos, establecimiento de enlaces y prevención de ciclos. Definir la estructura de comunicación lógica entre nodos, incluyendo la forma en que los dispositivos intercambian información y mantienen la continuidad de la topologia de manera autónoma. Implementar una propuesta conceptual del protocolo en dispositivos Android, mediante una aplicación demostrativa que permita la creación y reorganización dinámica de la red. Evaluar el comportamiento del protocolo mediante pruebas de laboratorio controladas, midiendo métricas de conectividad lógica, tiempo de reconfiguración ante fallos de nodo y continuidad del flujo de información en escenarios de topologia variable.

\seccion{Alcances}

Los alcances del proyecto definen los limites lógicos y compromisos técnicos para el desarrollo del protocolo distribuido, estableciendo de manera clara los componentes que serán materializados en el sistema experimental y los resultados esperados al finalizar la investigación.

Diseño formal del protocolo: Un documento técnico detallando la arquitectura lógica del protocolo, que incluye la especificación de los mecanismos de descubrimiento de nodos, diagramas de estado para el establecimiento de enlaces y tablas de decisión para la prevención de ciclos. Aplicación Android demostrativa: Un paquete de instalación (APK) funcional para dispositivos móviles que integra el motor del protocolo mediante interfaces Wi-Fi estándar, habilitando la formación de topologias dinámicas y un sistema de mensajeria en tiempo real. Validación en laboratorio: Un informe de resultados empiricos que consolida la evaluación de la conectividad inicial, el tiempo de reconexión topológica ante fallos y las métricas de continuidad del flujo de información en escenarios controlados.

\textbf{No entregables}

Modificaciones a la capa fisica o a la capa MAC del estándar IEEE 802.11. Mecanismos de calidad de servicio (QoS) o cifrado extremo a extremo. Despliegue en producción o compatibilidad con sistemas operativos distintos de Android.

\seccion{Limitaciones}

Tiempo limitado a 12 semanas según calendario académico semestral, restringiendo iteraciones de prueba.

El desarrollo está restringido a dispositivos con **Android 12 (API nivel 31)**, los cuales deben cumplir con los requisitos fisicos minimos del Compatibility Definition Document (CDD) de Google (Google, 2021): al menos 2 GB de memoria RAM (evitando la versión Go Edition para asegurar estabilidad), una resolución de pantalla minima de 426 dp x 320 dp y soporte mandatorio para Wi-Fi Direct (P2P), esencial para la formación de la red multisalto sin infraestructura. Acceso a funciones (SoftAP, Wi-Fi Direct) sujeto a restricciones de las APIs de Android por fabricante, dificultando la portabilidad. Al operar en capa de aplicación, el protocolo hereda la latencia y restricciones energéticas del SO, limitando los tiempos de respuesta.

\newpage
\capitulo{Marco teórico}

\seccion{Fundamentos conceptuales y tecnológicos}

El desarrollo del presente proyecto requiere una comprensión sólida de los fundamentos conceptuales y tecnológicos que sustentan el diseño de un protocolo distribuido para redes multisalto sobre Wi-Fi convencional. Este apartado tiene como finalidad proporcionar una visión integral de los conceptos relacionados con las redes inalámbricas, los protocolos de comunicación distribuida, las técnicas de enrutamiento y las herramientas tecnológicas empleadas en el desarrollo. A continuación, se presentan los elementos teóricos más relevantes, organizados en definiciones no técnicas y técnicas, con el propósito de contextualizar tanto la problemática abordada como las decisiones de diseño adoptadas en el protocolo propuesto.

Definiciones No Técnicas Se define como la relación entre fenómenos naturales peligrosos (terremotos, inundaciones, huracanes) y condiciones de vulnerabilidad fisica y socioeconómica de una comunidad (Maskrey, 1993, pp. 6--9). Para el proyecto, este concepto delimita el escenario de aplicación: entornos donde la infraestructura de comunicación convencional puede estar comprometida o destruida, y donde el protocolo propuesto debe operar de manera autónoma.

Resiliencia en Comunicaciones La resiliencia se entiende como la capacidad de un sistema para mantener o restaurar sus funciones esenciales frente a perturbaciones (Sterbenz et al., 2010). En este proyecto se materializa en la capacidad del protocolo para reorganizar la topologia ante la pérdida de nodos, garantizando la continuidad de la comunicación sin intervención del usuario.

Comunicación Descentralizada Se refiere a los mecanismos de intercambio de información en los que no existe un nodo central que coordine el tráfico; cada dispositivo toma decisiones autónomamente con información local (Conti \& Giordano, 2014). Este principio es el fundamento de diseño del protocolo propuesto, que permite a los nodos autoorganizarse sin depender de ningún coordinador externo.

Definiciones Técnicas Red de Área Local Inalámbrica (WLAN) y Estándar IEEE 802.11 Una WLAN emplea señales de radiofrecuencia para conectar dispositivos sin cableado; el estándar IEEE 802.11 (Wi-Fi) define los protocolos de las capas fisica y MAC (IEEE, 2020). En su modo convencional opera con un punto de acceso centralizado cuyo rango oscila entre 10 y 100 metros (Gupta \& Kumar, 2000), limitación que el protocolo propuesto busca superar operando en la capa de aplicación.

Red Ad Hoc y Redes Mesh Una red ad hoc es una red inalámbrica temporal y dinámica sin infraestructura preexistente, donde cada nodo actúa como cliente y como enrutador (Perkins, 2001). Las redes mesh evolucionan este concepto formando una malla con múltiples caminos alternativos que incrementa la robustez y la cobertura (Akyildiz et al., 2005), aunque su implementación comercial suele requerir hardware especializado que limita su uso en dispositivos móviles cotidianos (IEEE, 2020).

Una red multisalto permite la comunicación entre dispositivos fuera del rango directo mediante el reenvio de paquetes a través de nodos intermedios (Gupta \& Kumar, 2000). Este mecanismo es el núcleo del protocolo propuesto: permite que dispositivos Android con Wi-Fi convencional se comuniquen sin conexión directa, y la cantidad de saltos refleja la distancia topológica con implicaciones directas en latencia y confiabilidad.

Un protocolo distribuido es un conjunto de reglas que permite a múltiples nodos coordinarse y tomar decisiones autónomamente sin controlador centralizado (Tanenbaum \& Van Steen, 2007). En el proyecto, este componente habilita a los dispositivos Android para descubrir vecinos, establecer enlaces lógicos, prevenir ciclos y reorganizarse dinámicamente, convergiendo a un estado estable sin coordinación global.

Topologia de Red y Gestión Topológica La topologia describe cómo están interconectados los nodos; en redes dinámicas puede variar por movilidad, pérdida de enlaces o incorporación de nuevos dispositivos (Karl \& Willig, 2005). La gestión topológica engloba los mecanismos para mantener una representación coherente de la red y propagar cambios; en el protocolo propuesto este proceso ocurre de forma completamente distribuida.

Continuidad Lógica La continuidad lógica es la capacidad del protocolo para mantener el flujo de información incluso ante cambios topológicos, como la desconexión de un nodo o la degradación de un enlace (Karl \& Willig, 2005). Constituye uno de los objetivos centrales del protocolo y una de las métricas de evaluación principal en las pruebas experimentales, midiendo tiempo de reconexión y paquetes perdidos durante eventos de reconfiguración.

Sockets y Comunicación TCP/UDP en Android Un socket permite a una aplicación enviar y recibir datos en red; en Android se implementa sobre la API estándar de Java (java.net.Socket para TCP y java.net.DatagramSocket para UDP) sin acceder a niveles inferiores de la pila (Android Developers, 2024). El protocolo usa sockets TCP para intercambio de datos entre nodos con ruta establecida, y sockets UDP para mensajes de descubrimiento y señalización topológica, aprovechando su bajo overhead y compatibilidad con broadcast en la subred local (Tanenbaum \& Wetherall, 2011).

En definitiva, los conceptos y tecnologias revisados a lo largo de este apartado forman la base sobre la cual se sustenta el diseño del protocolo propuesto. Entender cómo funcionan las redes ad hoc y multisalto, qué estrategias existen para el enrutamiento y la prevención de ciclos, y qué herramientas ofrece la plataforma Android para la comunicación en red, permite tomar decisiones de diseño sólidas y evaluar con criterio el comportamiento del sistema bajo distintos escenarios de prueba.

\seccion{Estado del arte}

Ad hoc On-Demand Distance Vector (AODV) AODV calcula rutas bajo demanda difundiendo mensajes RREQ e instalando caminos de retorno RREP, usando secuencias para evitar ciclos (Perkins et al., 2003). Al exigir operación en la capa de red del SO, es inviable en Android no rooteado. El proyecto hereda su mitigación de fallos bajo demanda, trasladándola por completo a la capa de aplicación.

Optimized Link State Routing (OLSR) Protocolo proactivo que mapea rutas propagando continuamente el estado de los enlaces (Clausen \& Jacquet, 2003). Al igual que AODV, su dependencia estricta de la capa de red impide su ejecución en Android retail. De OLSR se rescata únicamente el concepto de mantener y actualizar topologias locales eficientemente.

Meshtastic implementa redes multisalto de largo alcance usando tecnologia LoRa para emergencias (Meshtastic Project, 2023). Su principal limitación es la dependencia de hardware especializado y una tasa de transferencia reducida (del orden de Kbps), lo que lo hace poco práctico para comunicación en tiempo real frente al Wi-Fi estándar.

Bridgefy y Briar Bridgefy es una plataforma de mensajeria descentralizada que emplea Bluetooth y Wi-Fi para comunicación sin internet, pero investigadores de la ETH Zúrich identificaron vulnerabilidades serias en su protocolo, incluyendo suplantación de identidad y ausencia de confidencialidad (Albrecht et al., 2021); adicionalmente, su código cerrado dificulta su uso académico. Briar, por su parte, es una aplicación de código abierto con cifrado extremo a extremo sobre Bluetooth y Wi-Fi (Briar Project, 2021), pero su enrutamiento multisalto está restringido a contactos previamente conocidos, limitando la cobertura en redes con nodos desconocidos. Ambas confirman la viabilidad técnica de la comunicación descentralizada sobre hardware móvil estándar, pero ninguna aborda la reconfiguración topológica distribuida ante cambios dinámicos.

BitChat es un protocolo de mensajeria descentralizado y cifrado de extremo a extremo que utiliza una arquitectura hibrida de red en malla sobre Bluetooth Low Energy (BLE) y el protocolo Nostr para el enrutamiento global (Biagioni, 2025). A diferencia de las soluciones centralizadas, opera sin identificadores persistentes ni servidores de coordinación, permitiendo la comunicación en entornos sin acceso a internet o redes celulares. Aunque implementa un sistema multihop de hasta 7 saltos, su dependencia del estándar Bluetooth limita el ancho de banda efectivo y la profundidad de la red, lo que lo posiciona como una solución de baja tasa de bits frente a las capacidades teóricas del Wi-Fi.

Análisis comparativo Protocolo / App Hardware AODV / OLSR Wi-Fi Estándar Si (Capa Red) Radios LoRa No Bridgefy / Briar Bluetooth/Wi-Fi No (Capa App) Bluetooth Mesh No (Capa App) Propuesta actual Wi-Fi Estándar No (Capa App)

\seccion{Marco legal}

El desarrollo del proyecto involucra software, datos de usuario y comunicaciones inalámbricas, por lo que debe alinearse con el marco normativo colombiano vigente en materia de protección de datos personales, TIC y derechos de autor.

Ley 1581 de 2012 --- Protección de datos personales Regula el tratamiento de datos personales en Colombia, exigiendo estándares de confi- dencialidad (Congreso de la República de Colombia, 2012). Para el proyecto, obliga a que los mensajes ajenos no sean retenidos prolongadamente por los nodos intermediarios, descartándolos de la memoria base del dispositivo tan pronto concluya el salto estructural.

Ley 1341 de 2009 --- Ley de TIC Esta ley establece los principios del sector TIC en Colombia y promueve el acceso universal a las comunicaciones, la investigación tecnológica y la neutralidad tecnológica (Congreso de la República de Colombia, 2009). El proyecto se alinea con esta ley en tanto el protocolo propuesto busca extender la conectividad en escenarios donde los servicios convencionales no están disponibles.

\newpage
\capitulo{Diseño metodológico}

Para el desarrollo del presente proyecto se implementará una metodologia incremental. Este enfoque de ingenieria de software se basa en la construcción del sistema a través de entregas parciales y sucesivas, denominadas incrementos. En cada incremento se realiza un ciclo que abarca el análisis, diseño, implementación y pruebas de un conjunto especifico de requisitos funcionales, lo que permite validar progresivamente el protocolo propuesto. A diferencia de un modelo lineal tradicional, el desarrollo iterativo facilita la retroalimentación continua y la detección temprana de errores conceptuales o técnicos. Esta caracteristica resulta fundamental para el diseño de un protocolo sobre capas de abstracción no nativas (capa de aplicación en Android), donde limitaciones imprevistas de las APIs o comportamientos anómalos del entorno de red requieren ajustes dinámicos de diseño sin afectar todo el ciclo de vida del proyecto.

El equipo de trabajo está conformado por integrantes que asumirán roles complementarios durante las iteraciones, siguiendo lineamientos inspirados en marcos de trabajo ágiles adaptados al desarrollo académico: Equipo de Desarrollo (Estudiantes proyectantes): • Juan Carlos Clavijo: Responsable del diseño lógico del protocolo, la arquitectura concurrente, el desarrollo de los algoritmos de enrutamiento (tablas de salto) y la lógica de detección de fallos mediante heartbeats. • Brandon Stiven Ganzo: Encargado de la implementación de la capa de comunicación de bajo nivel (UDP broadcast), el túnel de control TCP, el desarrollo del subsistema de mensajeria y el diseño de la interfaz gráfica de usuario. Director del Proyecto: Cumple el rol de orientador metodológico y técnico. Su responsabilidad es velar por el rigor del proceso, validar la coherencia conceptual del diseño con la teoria de redes distribuidas y evaluar los incrementos desarrollados, asegurando que cada entrega aporte al cumplimiento de los objetivos especificos.

Fases del proyecto y desarrollo iterativo El proyecto se organizará en cuatro iteraciones o incrementos principales correspondientes al desarrollo del protocolo y su aplicación en el sistema demostrativo. Cada fase culmina con un entregable verificable.

Incremento 1: Análisis y diseño de la arquitectura base Este incremento establece los cimientos lógicos del sistema. Incluye la revisión de los estándares técnicos (Wi-Fi Manager, Wi-Fi Direct) y protocolos de red ad hoc de referencia (AODV, OLSR), definiendo cómo se aplicarán las reglas de enrutamiento distribuidas en un entorno sin infraestructura. Actividades: Identificación de requerimientos, análisis de restricciones en APIs de Android (desde Android 12 en adelante), diseño lógico de las estructuras de datos de enrutamiento y concepción arquitectónica del sistema operativo concurrente. Entregables: Documento de requerimientos priorizado, diagramas de clases y de secuencia del protocolo de comunicación base.

Incremento 2: Descubrimiento de nodos y topologia inicial Se procede con la implementación de los servicios de anuncio y la formación de una red local autónoma mediante la asignación de roles de acceso (SoftAP) y puntos cliente. Actividades: Desarrollo del módulo de anuncio por broadcast UDP, programación de los sockets TCP para el establecimiento de enlaces seguros y creación de la interfaz (UI) básica de inicio de sistema en Android. Entregables: Aplicación capaz de identificar vecinos en el mismo espectro de cobertura Wi-Fi y establecer un enlace de una via entre dos nodos sin intervención externa.

Incremento 3: Enrutamiento y mensajeria en red multisalto Este incremento corresponde al núcleo del proyecto. Sobre los nodos enlazados, se diseñan e implementan las tablas de estados de los vecinos (routing tables) que posibilitan extender la comunicación a nivel escalonado. Se integran mensajes de control inter-nodos (señalizadores) y mensajes de carga útil (texto o chat de usuario). Actividades: Desarrollo del algoritmo de propagación de rutas basándose en TTL y secuencias numéricas, control de ciclos infinitos, e implementación del intercambio estructurado de mensajes entre nodos no directos. Entregables: Software beta con capacidad de intercambio de mensajes bidireccional abarcando una cadena minima de 3 dispositivos intermedios.

Incremento 4: Reconfiguración topológica, validación y escalamiento Fase final de análisis dinámico donde ocurren las disrupciones para validar las facultades resilientes. Se incorporan las funciones lógicas de detección de desconexión abrupta por medio de ”heartbeatsregulares o fallos de envios. Actividades: Programación del procedimiento de purga de rutas rotas tras intervalos temporizados, programación de re-generación e incrustación de métricas automáticas, y pruebas empiricas experimentales completas. Entregables: Versión estable con reconexión reactiva demostrable ante apagados sistemáticos, recolección depurada de datos empiricos e informe con los manuales/documentación. Para propósitos de organización, la Tabla 3.1 detalla las acciones condensadas de estos procesos en sus respectivos incrementos funcionales.

Incremento Tareas principales Levantamiento de requerimientos, análisis de restricciones en Android, diseño de arquitectura concurrente.

Programación UDP broadcast, establecimiento de Conectividad túnel de transmisión TCP-IP en dispositivos en Base corto rango. 3. Subsistema Programación lógica del intercambio de tablas de (routing), sistema de contención anti-ciclos e enrutamiento interfaz de mensajeria (chat). 4. Estabilidad Lógica de latidos, sistema de reconstrucción y Pruebas automática tras la caida de puentes e inyección de datos experimentales. 1. Análisis y Diseño

Requisitos documentados, diagramas UML / de Secuencia. Módulo de escaneo emparejado entre 2 nodos. Red conformable por $\geq$ 3 saltos simultáneos y funcionales. Versión final, registros empiricos, manual y documentaciones finales.

Técnicas e instrumentos de validación experimental Dado el enfoque puramente tecno-experimental, la validación del trabajo reposa sobre el desempeño real del producto derivado de la programación en laboratorio y no de análisis sociodemográficos u observación participativa teórica. Los ensayos se desarrollarán creando redes preconfiguradas y redes emergentes (espontáneas) empleando múltiples terminales móviles controlados. Se evaluará el cumplimiento de cada requisito funcional establecido en el incremento preliminar. Las variables observadas consistirán en: Tiempo de reconexión adaptativo: Fracción en milisegundos que la red invertirá hasta encontrar y propiciar un nuevo recorrido efectivo cuando el enlace principal se interrumpe artificialmente. Estabilidad operativa y profundidad: Verificación del máximo de escalonamiento estructural (saltos) permitidos en función del ancho de banda y latencia originada por la retransmisión computarizada de la pasarela de Android.

